\documentclass[12pt]{article}
\usepackage[T1]{fontenc}
\usepackage[utf8]{inputenc}
\usepackage{lmodern}
\usepackage{textcomp}
\usepackage{enumitem}
\usepackage{geometry}
\geometry{margin=1in}
\begin{document}
\begin{center}
Answer Key - Psychology - Graduate Level Assessment 22 \
\small (Answers are ordered exactly as in the question paper.)
\end{center}

\section*{Multiple Choice}
\begin{enumerate}[label=\arabic*.]
\item \textbf{Answer:} A
\\ \textit{Options:} A) Does not support self-efficacy for change; B) Empathy through reflective listening; C) Roll with resistance; D) Avoid direct confrontation
\\ \textit{Note:} The correct answer is "Does not support self-efficacy for change" because Motivational Interviewing fundamentally emphasizes enhancing a patient's self-efficacy, which is critical for fostering commitment and facilitating decision-making in the change process.
\item \textbf{Answer:} A
\\ \textit{Options:} A) lack of debriefing after the experiment; B) lack of compensation for participants; C) violation of anonymity; D) nonrandom sampling procedures; E) physical harm to participants
\\ \textit{Note:} The correct answer pertains to the fact that participants in Milgram's obedience study experienced significant psychological stress and were not adequately debriefed afterward, raising serious ethical concerns about informed consent and the duty of researchers to ensure the well-being of their subjects.
\item \textbf{Answer:} A
\\ \textit{Options:} A) orange; B) pink; C) brown; D) red; E) black
\\ \textit{Note:} The correct answer is orange because of the opponent-process theory of color vision, which posits that staring at a yellow stimulus causes the visual system to fatigue the yellow receptors, leading to the perception of its complementary color, orange, when the stimulus is removed.
\item \textbf{Answer:} B
\\ \textit{Options:} A) Homeostatic regulation; B) Hierarchy of needs; C) Self-determination theory; D) Intrinsic motivation; E) Maslow's pyramid of desire
\\ \textit{Note:} The correct answer is the "Hierarchy of needs" because Maslow's theory categorizes human motivations into a five-tier model, prioritizing basic physiological needs at the base as essential for survival, followed by safety, love and belonging, esteem, and self-actualization, thereby illustrating that some motives are indeed more imperative than others.
\item \textbf{Answer:} A
\\ \textit{Options:} A) Mean: 14 years, Median: 12.5, Mode: 13, Range: 8; B) Mean: 13 years, Median: 14, Mode: 12.5, Range: 17; C) Mean: 13.5 years, Median: 12.5, Mode: 14, Range: 8; D) Mean: 13 years, Median: 13, Mode: 8, Range: 9; E) Mean: 12.5 years, Median: 12, Mode: 11, Range: 10
\\ \textit{Note:} The correct answer is based on the calculations of the mean (average of the ages), median (middle value when the ages are ordered), mode (most frequently occurring age), and range (difference between the highest and lowest ages), which accurately reflect the measures of central tendency and variability for the provided data set.
\end{enumerate}

\section*{True / False}
\begin{enumerate}[label=\arabic*.]
\setcounter{enumi}{5}
\item \textbf{Answer:} False
\\ \textit{Options:} A) True; B) False
\\ \textit{Note:} Primary prevention encompasses both individual-level interventions and population-wide approaches aimed at reducing the incidence of disorders or health issues before they occur.
\item \textbf{Answer:} True
\\ \textit{Options:} A) True; B) False
\\ \textit{Note:} The DSM-5 diagnostic criteria for problem gambling emphasize criteria such as preoccupation, tolerance, and withdrawal symptoms, but do not specifically list loss of control as a defining feature, making the statement true.
\item \textbf{Answer:} False
\\ \textit{Options:} A) True; B) False
\\ \textit{Note:} Affective commitment, while important for employee engagement and job satisfaction, does not directly correlate with productivity levels, as productivity is influenced by various other factors such as skill level, work environment, and task complexity.
\item \textbf{Answer:} False
\\ \textit{Options:} A) True; B) False
\\ \textit{Note:} The letter 'c' in the word 'cat' functions as a grapheme representing a single sound, specifically the /k/ sound, rather than a combination of sounds.
\item \textbf{Answer:} True
\\ \textit{Options:} A) True; B) False
\\ \textit{Note:} Chemotherapy is considered a somatic therapy because it directly targets the physiological aspects of cancer by using chemical agents to kill or inhibit the growth of cancer cells, thereby eliciting specific bodily responses.
\end{enumerate}

\section*{Long Form Response}
\begin{enumerate}[label=\arabic*.]
\setcounter{enumi}{10}
\item \textbf{Answer:} Echoic memory serves as a temporary storage system for auditory information, characterized by its brief duration of approximately one to ten seconds. This type of memory allows individuals to retain sounds long enough for processing and comprehension. Unlike a tape recorder, which captures and reproduces sound with complete fidelity, echoic memory involves a transformation of auditory stimuli through the cochlea and possibly other cognitive processes. This means that by the time auditory information enters echoic memory, it has already undergone an initial filtering and interpretation, impacting how it is later retrieved and utilized. Thus, while both systems store information, echoic memory is fundamentally more complex and influenced by perceptual and cognitive mechanisms.
\\ \textit{Note:} correct because it accurately describes echoic memory as a temporary auditory storage system with a brief duration, emphasizes its role in processing sounds rather than merely recording them like a tape recorder, and highlights the cognitive transformations involved in auditory perception.
\item \textbf{Answer:} Cognitive dissonance refers to the psychological discomfort experienced when an individual holds two contradictory beliefs or makes conflicting decisions. In Delia's decision-making process between attending Harvard University and Yale University, she may experience cognitive dissonance if she values both institutions highly but must ultimately choose one, leading to feelings of regret or uncertainty about her decision. This dissonance can manifest as anxiety or self-doubt, particularly if she perceives that the alternative choice (the other university) might offer opportunities she would miss. To resolve this dissonance, Delia might seek to reinforce her choice by focusing on the positive aspects of Harvard, such as its renowned programs or networking opportunities, while minimizing the perceived benefits of Yale. This psychological mechanism underscores the complexities inherent in decision-making, particularly in high-stakes situations where personal and academic aspirations are at play.
\\ \textit{Note:} correct because it accurately describes cognitive dissonance as the psychological conflict arising from choosing between two valued options-in this case, Harvard and Yale-resulting in potential regret and uncertainty in Delia's decision-making process.
\end{enumerate}

\vfill
\noindent\textit{End of Answer Key}
\end{document}